\section{Estrutura do pitch}
\label{sec:estrutura}

A estrutura segue o encadeamento problema $\rightarrow$ solução $\rightarrow$ mercado $\rightarrow$ negócio $\rightarrow$ evidência $\rightarrow$ futuro. Cada bloco responde a uma pergunta que o público faz, de modo natural, ao ouvir uma proposta de startup. A Tabela~\ref{tab:estrutura} relaciona os blocos, a pergunta respondida, a mensagem principal e a lâmina correspondente.

\begin{table}[htb]
\centering
\caption{Estrutura do pitch: blocos, perguntas e mensagens}
\label{tab:estrutura}
\fontsizetable
\begin{tabularx}{\textwidth}{L{0.9cm} L{2.7cm} L{3.7cm} Y}
\toprule
\textbf{Slide} & \textbf{Bloco} & \textbf{Pergunta que responde} & \textbf{Mensagem principal}\\
\midrule
1 & Abertura & Quem somos? & PetCuida: cuidar antes que vire emergência.\\
2 & História & Por que isso importa? & Um tutor sem acesso a prevenção só chega ao veterinário na emergência.\\
3 & Problema & Qual é a dor, e quão grande é? & Há milhões de lares com cães e gatos, mas o cuidado preventivo é caro e fragmentado.\\
4 & Solução & O que oferecemos? & Um aplicativo que organiza a saúde do animal e conecta o tutor a clínicas parceiras.\\
5 & Público-alvo & Para quem? & Tutores de baixa renda (núcleo) e clínicas e ONGs (o outro lado da plataforma).\\
6 & Proposta de valor & Por que escolher? & Prevenção gratuita para o tutor; clientes e organização para a clínica.\\
7 & Modelo de negócio & Como ganhamos dinheiro? & Núcleo gratuito; receita de comissão, assinatura opcional, plano para clínicas e convênios.\\
8 & Diferenciais & Por que não outro? & Não é plano de saúde nem apenas agenda: é prevenção com rede solidária e preço social.\\
9 & Tração & O que já foi feito? & Protótipo funcional, cinco sessões de teste e economia unitária modelada.\\
10 & Próximos passos & O que vem agora? & Piloto em Belo Horizonte, teste de preço, formalização e segurança de dados.\\
11 & Equipe e pedido & Quem executa e o que precisamos? & Quatro integrantes com papéis definidos; buscamos clínicas parceiras, ONGs e mentoria.\\
\bottomrule
\end{tabularx}
\fonteautores
\end{table}

Os oito tópicos exigidos pela orientação estão integralmente cobertos (lâminas 3 a 10). A abertura, a história e o fechamento foram acrescentados para dar ao pitch início, arco narrativo e pedido explícito, pois uma apresentação sem pedido deixa o público sem saber qual ação tomar.
